% =====================================================================================
% Note pour Hugo Rapior : relevé des livrables publics du Lab Cyber et trajectoire 2030.
% Etabli le 2 septembre 2026, uniquement sur SOURCES PUBLIQUES.
%
% PRECAUTION DE TON, ET ELLE COMMANDE TOUT LE DOCUMENT. Hugo cosigne les deux rapports LUCY
% et anime l'atelier du 24 novembre : il connait son Lab mieux que moi. Ce document ne lui
% apprend donc pas ce que fait le Lab, il lui dit CE QU'UN TIERS EN VOIT, ce qui est une
% information differente et qu'il n'a aucune raison d'avoir. D'ou le cadrage annonce des la
% premiere phrase, et le refus de donner des conseils de staffing sous forme d'affirmations.
%
% DEUX CONTRAINTES DE FORME :
%   - aucun tiret cadratin. Le babel francais en met un en puce d'itemize par defaut ;
%   - tout ce qui n'est pas corrobore par une source publique est signale comme tel, au meme
%     titre que le memoire le fait pour Hackmageddon et que sa note interne le fait pour PANAME.
% =====================================================================================
\documentclass[11pt,a4paper]{article}

\usepackage[utf8]{inputenc}
\usepackage[T1]{fontenc}
\usepackage{lmodern}
\usepackage[french]{babel}
\usepackage[top=1.5cm,bottom=1.4cm,left=1.85cm,right=1.85cm]{geometry}
\usepackage{booktabs}
\usepackage{longtable}
\usepackage{array}
\usepackage{enumitem}
\usepackage{xcolor}
\usepackage{eurosym}          % \euro : absent du codage T1, il faut le paquet
\usepackage{titlesec}

% Charte Nexialog : memes roles que style_nexialog.py, aucun hexadecimal ailleurs qu'ici.
\definecolor{encre}{HTML}{1B1E30}
\definecolor{accent}{HTML}{A6002E}
\definecolor{bleu}{HTML}{204993}
\definecolor{teal}{HTML}{007D78}
\definecolor{grisclair}{HTML}{F2F2F2}

\color{encre}
\titleformat{\section}{\normalfont\large\bfseries\color{bleu}}{\thesection.\ }{0.4em}{}
\titlespacing{\section}{0pt}{9pt}{3pt}
\titleformat{\subsection}{\normalfont\bfseries}{}{0pt}{}
\titlespacing{\subsection}{0pt}{9pt}{2pt}
% Puce explicite : le babel francais mettrait un tiret cadratin.
\setlist[itemize]{label=\textcolor{accent}{\textbullet},itemsep=2.5pt,topsep=3pt,
                  leftmargin=1.15em,parsep=0pt}
\renewcommand{\arraystretch}{1.22}
\newcolumntype{L}[1]{>{\raggedright\arraybackslash}p{#1}}

\newcommand{\etat}[2]{\textcolor{#1}{\textbf{\footnotesize #2}}}
\newcommand{\encadre}[2]{%
  \vspace{3pt}\noindent\colorbox{grisclair}{%
  \begin{minipage}{\dimexpr\textwidth-2\fboxsep}
  \small\textbf{#1} #2
  \end{minipage}}\vspace{4pt}}

\pagestyle{plain}

\begin{document}
\small

{\Large\bfseries\color{bleu} Lab Cyber : les livrables visibles, et une trajectoire 2030}\\[3pt]
{\color{accent}\rule{\textwidth}{1.1pt}}\\[3pt]
{\footnotesize Kélian Kaddouri \hfill Pour Hugo Rapior \textbullet{} 2 septembre 2026}

\vspace{6pt}

\noindent\textbf{Ce que ce document est, et ce qu'il n'est pas.} Tu connais le Lab mieux que
moi, tu cosignes les deux rapports LUCY et tu animes l'atelier de novembre. Ce relevé ne te dit
donc pas ce que fait le Lab : il te dit \textbf{ce qu'un tiers en voit} quand il ne dispose que
des sources publiques, ce qui est une information différente. Je l'ai fait pour savoir où mon
mémoire se branche, et il en est sorti un constat qui me semble valoir le détour. Tout ce qui
n'est pas corroboré par une source publique est signalé comme tel.

\section{Les cinq axes annoncés, et les livrables qui les portent}

La page du pôle R\&D énumère cinq axes pour le Lab Cyber. Confrontés aux publications
réellement en ligne, trois sont portés par un livrable, trois ne le sont pas.

\begin{small}
\begin{longtable}{@{}L{5.2cm} L{8.6cm} l@{}}
\toprule
\textbf{Axe annoncé} & \textbf{Livrable public qui le porte} & \textbf{État} \\
\midrule
\endfirsthead
\toprule
\textbf{Axe annoncé} & \textbf{Livrable public} & \textbf{État} \\
\midrule
\endhead
\bottomrule
\endfoot
Modélisation actuarielle du risque cyber
& Atelier \og{}Scénarios stochastiques, données textuelles et IA pour une meilleure
  quantification du risque cyber\fg{}, journée 100 \% Actuaires, 24 novembre 2025
& \etat{teal}{porté} \\
\addlinespace
Marché de la cyberassurance, hors liste des cinq axes mais c'est le livrable cyber le plus
  visible
& Rapports LUCY 2025 et LUCY 2026, analyse du marché français
& \etat{teal}{porté} \\
\addlinespace
Audit de confidentialité des modèles d'IA
& Livre blanc \og{}Fondations et gouvernance pour une IA de confiance\fg{} avec Finance
  Innovation, et l'article du 10 juillet 2025 sur le GPAI Code of Practice
& \etat{teal}{porté} \\
\addlinespace
Conformité DORA et NIS2 : cartographie ICT, KPI et KRI
& Aucune publication dédiée trouvée. DORA n'apparaît que mentionnée, dans LUCY, comme
  réglementation applicable aux petites structures
& \etat{accent}{non porté} \\
\addlinespace
Qualification et scoring des tiers ICT critiques
& Aucune publication trouvée
& \etat{accent}{non porté} \\
\addlinespace
\textbf{Calcul du SCR cyber (assurance) et RWA cyber (banque)}
& \textbf{Aucune publication.} LUCY 2025 ne mentionne pas le SCR une seule fois, et l'atelier
  de novembre ne mentionne ni DORA ni SCR
& \etat{accent}{non porté} \\
\end{longtable}
\end{small}

\encadre{Le constat, en une phrase.}{Le Lab annonce le calcul du SCR cyber et du RWA cyber, et
c'est le seul de ses axes qui ne soit porté par aucune publication. C'est aussi, à la virgule
près, le sujet de mon mémoire.}

\subsection{Le détail des trois livrables, puisque c'est ce qui fait foi}

\begin{itemize}
\item \textbf{Les rapports LUCY.} Signés Hugo Rapior et Franck Osée Dountio Zaboue, assis sur
      14\,124 polices et 448 sinistres déclarés auprès d'Aon, Marsh et WTW. Ce qu'ils
      établissent : un ratio de sinistralité de 17 \% en 2024 contre 12 \% en 2023, un nombre de
      sinistres en baisse de 27 \% mais des montants indemnisés en hausse de 43 \%, le retour des
      sinistres XXL avec 21 M\euro{} sur les incidents majeurs, des tarifs en baisse de 19,8 \%
      sur les grandes entreprises et de 8,7 \% sur les ETI, des capacités en hausse de 23 \%, et
      des franchises en baisse de 15 \% pour les grands comptes contre une hausse de 231 \% pour
      les PME. C'est un livrable de \textbf{marché}.
\item \textbf{L'atelier du 24 novembre 2025.} Avec Franck Dountio, Chadi Hanna pour Dattak et
      Caroline Hillairet. Trois méthodes annoncées : générateurs de scénarios stochastiques face
      aux pandémies cyber, processus de Poisson \textbf{auto-excités} intégrant contagion et
      divulgation de vulnérabilités, et exploitation de données textuelles par NLP et grands
      modèles de langage jusqu'aux transformeurs pour la sévérité. C'est un livrable de
      \textbf{méthode}.
\item \textbf{Le livre blanc IA de confiance}, avec Finance Innovation, porté par l'équipe FCC :
      qualité de la donnée, dispositifs de gouvernance et de contrôle, cas d'usage sur les trois
      lignes de défense, et un parangonnage des solutions d'IA déployées en banque et en
      assurance dont la fraude et la cybersécurité. Complété par l'article sur le GPAI Code of
      Practice, adossé à l'article 53 de l'IA Act.
\end{itemize}

\section{Trois lignes que les sources publiques ne corroborent pas}

Ce n'est pas un reproche, et il faut le dire d'emblée : une intervention orale, une veille ou
une contribution informelle ne laissent pas de trace publique. C'est une information sur la
\textbf{visibilité externe} du Lab, et elle est utile parce qu'un interlocuteur qui vérifie
trouvera exactement ce que j'ai trouvé.

\begin{itemize}
\item \textbf{PANAME.} Projet public piloté par l'ANSSI, la CNIL, le PEReN et Inria. Le Lab a
      bien l'axe \og{}audit de confidentialité des modèles d'IA\fg{}, mais aucune publication ne
      permet d'attribuer un livrable Nexialog sous ce label. La note interne le signale déjà,
      et je le confirme.
\item \textbf{La présentation au European Congress of Actuaries.} Le congrès a lieu à Paris les
      18 et 19 juin 2026 et son appel à intervenants comportait bien \og{}Role of actuaries in
      modelling Cyber Risk\fg{} en risques émergents. Aucune intervention Nexialog n'apparaît
      dans les sources consultées.
\item \textbf{Le risque géopolitique.} Il figure dans la présentation générale du cabinet parmi
      les risques transverses sur lesquels il innove, mais aucune publication dédiée n'est en
      ligne.
\end{itemize}

\section{Où mon mémoire se branche, et le point qu'il faut que je te dise}

La méthode publiée du Lab repose sur des processus de Poisson \textbf{auto-excités}, donc un
Hawkes. Mon mémoire l'écarte au profit d'une cascade dirigée entre les cinq piliers. Présenté
ainsi c'est un désaccord, et ce serait une mauvaise façon d'arriver.

\encadre{Ce qui réconcilie les deux, et c'est un résultat plutôt qu'un arrangement.}{Une cascade
\emph{sans aucune} auto-excitation reproduit un ratio de branchement de 0,480 au pas journalier
contre 0,551 mesuré, et retrouve même la demi-vie. Les deux familles sont donc
\textbf{observationnellement équivalentes} à la résolution disponible : le rejet du Hawkes se
reformule en équivalence suivie d'un choix de parcimonie interprétative. Ce résultat ne
contredit pas la ligne du Lab, il la \textbf{borne}, et il coupe dans les deux sens puisqu'il
interdit aussi de revendiquer qu'une des deux paramétrisations est la bonne \emph{parce que} la
donnée le dirait. Ce qui se défend est la traçabilité et la déclaration des limites, pas la
supériorité du modèle.}

La complémentarité est nette sur les objets. Le Lab publie le \textbf{marché} (LUCY) et la
\textbf{méthode} (générateurs, NLP). Mon mémoire publie le \textbf{capital} : besoin en pilier 2
au titre de la non-conformité DORA, de 6\,049 à 20\,188 M\euro{} selon l'état de conformité, soit
un facteur 3,34, avec 1\,980 nombres publiés dont chacun est rattaché au script qui le calcule.
Trois briques différentes de la même offre, et la troisième est celle qui n'a pas encore de
livrable.

\section{La trajectoire 2026-2030, telle que le calendrier la contraint}

\begin{small}
\begin{longtable}{@{}L{2.5cm} L{11.4cm}@{}}
\toprule
\textbf{Échéance} & \textbf{Ce que le jalon crée comme demande} \\
\midrule
\endfirsthead
\toprule
\textbf{Échéance} & \textbf{Ce que le jalon crée comme demande} \\
\midrule
\endhead
\bottomrule
\endfoot
30 avril 2026
& Registre d'information, date de référence au 31 décembre 2025. À l'exercice à blanc de 2024,
  \textbf{6,5 \% seulement} des près de mille entités testées passaient les 116 contrôles
  qualité. C'est ce chiffre qui fabrique du mandat aujourd'hui, et c'est autant de la qualité de
  donnée que de la conformité. \\
\addlinespace
2 août 2026
& Obligations de l'IA Act sur les systèmes à haut risque. Les modèles de scoring, de détection
  de fraude et les agents conversationnels deviennent des actifs à cartographier et à auditer :
  l'axe gouvernance de l'IA cesse d'être une veille. \\
\addlinespace
2026-2027
& Premier cycle de supervision directe des \textbf{19 prestataires TIC critiques} désignés le
  18 novembre 2025, dont AWS, Azure, Google Cloud, Bloomberg, LSEG, IBM, TCS et Orange.
  \textbf{C'est le jalon le plus exploitable du lot} : une liste nominative de dix-neuf points de
  concentration transforme l'accumulation par cause commune en exercice mesurable sur des tiers
  identifiés, au lieu d'un facteur posé. \\
\addlinespace
17 janvier 2027
& DORA pleinement applicable aux petites entités. Élargissement de l'assiette vers des acteurs
  sans direction des risques constituée : le besoin bascule vers de l'outillage et du modèle
  standardisé plutôt que du sur-mesure. \\
\addlinespace
17 janvier 2028
& Rapport de réexamen de DORA par la Commission, avec proposition législative possible.
  \textbf{C'est la seule fenêtre où une méthode publiée peut peser sur le cadre} plutôt que le
  subir, et elle se prépare deux ans à l'avance. \\
\addlinespace
2026-2030
& Stratégie nationale de cybersécurité. Le directeur général de l'ANSSI a fixé 2030 comme
  horizon d'une possible guerre cyber de haute intensité, avec risque de saturation des capacités
  de défense. Une menace corrélée et massive est le régime où un modèle de cascade et
  d'accumulation cesse d'être un raffinement. \\
\end{longtable}
\end{small}

\section{L'offre qui en découle, en quatre briques}

\begin{itemize}
\item \textbf{Conformité outillée, aujourd'hui.} Registre d'information, cartographie ICT, KPI et
      KRI, scoring des tiers critiques. Le mandat existe et le taux de 6,5 \% le prouve. Mais
      c'est la brique qui se \textbf{banalise le plus vite} : dès que l'exercice se répète, il
      passe à l'outil et le prix baisse. Elle finance le reste, elle ne fait pas l'identité du
      Lab.
\item \textbf{Quantification du capital, 2026-2027.} SCR cyber, RWA cyber, besoin en pilier 2 au
      titre de la non-conformité. C'est l'axe annoncé et non publié, celui qui distingue un
      cabinet d'actuariat d'un cabinet de cybersécurité, et celui qu'aucun outil ne remplace
      parce qu'il exige un modèle défendable devant un régulateur.
\item \textbf{Accumulation et transfert, 2027-2028.} La suite naturelle de LUCY et des dix-neuf
      prestataires critiques : accumulation par cause commune, dépendance de queue, structuration
      de la cession. Le pont entre le livrable de marché déjà publié et le modèle de capital, et
      le premier levier qu'un client actionne avant le capital.
\item \textbf{Audit de modèles, cyber et IA confondus, 2028-2030.} Les modèles d'IA deviennent
      des actifs TIC au sens du pilier 1 et l'ANSSI construit des schémas de certification. La
      compétence qui se vend alors n'est ni le cyber ni l'IA, c'est la \textbf{validation de
      seconde ligne} appliquée aux deux, avec la traçabilité comme livrable.
\end{itemize}

\encadre{Le fil qui relie les quatre.}{De 2026 à 2030, la demande passe de la
\textbf{documentation} à la \textbf{preuve} : registres vérifiés, tests pilotés par la menace,
audits de modèles effectifs. Savoir produire une chaîne de calcul dont chaque nombre se rattache
au script qui le produit, c'est vendre exactement cela.}

\section{Le risque qui peut faire échouer la brique 2, et il porte sur mon propre travail}

Toute la quantification repose sur une demande qui n'est pas \emph{exigée} : il n'existe aucun
module DORA en formule standard, donc le besoin de capital cyber reste un exercice d'ORSA, c'est
à dire un récit que l'entité choisit de produire. Si le réexamen de janvier 2028 ne crée pas
d'exigence chiffrée, la quantification restera une prestation de conviction et non de
conformité, avec un marché plus petit et des cycles de vente plus longs. C'est la limite que mon
mémoire publie lui-même, et c'est aussi la raison pour laquelle le jalon de 2028 mérite d'être
préparé plutôt qu'attendu.

Deux réserves secondaires, du même ordre. La \textbf{donnée} reste le facteur limitant : rien
dans le calendrier ne crée de base publique codant l'ordre des défaillances entre domaines de
contrôle, et le seul chemin identifié est la notification d'incident majeur sous DORA, qui est
exactement la recommandation de reporting que mon mémoire formule. Et la
\textbf{différenciation méthodologique} est difficile à défendre, pour la raison d'équivalence
observationnelle donnée plus haut.

\section{Trois observations sur l'équipe, sous forme de questions}

Je les pose plutôt que je ne les affirme : tu vois le carnet de commandes et le plan de charge
que ce relevé ne voit pas.

\begin{itemize}
\item Ce que les traces publiques laissent voir du Lab est \textbf{mince} : un responsable de
      programme, des stagiaires ENSAE de fin d'études, et un recrutement en cours de consultant
      leader R\&D cyber. Est-ce que la brique 2 attend ce recrutement, ou est-ce qu'elle peut
      démarrer avant ?
\item Le canal de recrutement visible est \textbf{unique}, l'ENSAE et le master SFA. C'est
      efficace, et c'est fragile.
\item Le pôle mène des \textbf{thèses CIFRE} avec Dauphine et le LEO sur la finance durable et
      le climat, et \textbf{aucune sur le cyber} alors que c'est l'un des trois labs. Une CIFRE
      cyber achèterait trois ans de recherche financée et un titre publiable, ce qu'aucune
      mission ne produit. Est-ce que le sujet a déjà été instruit ?
\end{itemize}

\section{Ce que je propose}

\begin{itemize}
\item \textbf{Publier la brique 2.} J'ai la chaîne de calcul, sa traçabilité et ses limites
      déclarées. La question n'est pas de savoir si je peux contribuer au Lab, c'est de savoir
      sous quelle forme on publie : note méthodologique interne réutilisable en avant-vente,
      article, ou les deux. La seconde forme est celle qui achète la légitimité externe qui
      manque à cet axe.
\item \textbf{Prendre les dix-neuf prestataires critiques comme prolongement.} C'est le
      prolongement le plus direct de mon pilier 4, et il tombe dans la fenêtre de supervision
      2026-2027.
\item \textbf{Préparer le réexamen de 2028} avec l'Institut des Actuaires, ce qui ne se fait pas
      seul et se prépare tôt.
\end{itemize}

\vspace{8pt}
\noindent{\color{accent}\rule{\textwidth}{0.6pt}}\\[2pt]
{\footnotesize\itshape Relevé du 2 septembre 2026 sur sources publiques : pages Nexialog (pôle
R\&D, publications, LUCY 2025 et 2026, atelier 100 \% Actuaires, livre blanc IA de confiance,
GPAI Code of Practice), ESA, cyber.gouv.fr, EUR-Lex, Actuarial Association of Europe. Les
chiffres du mémoire sont ceux du dépôt au 2 septembre 2026. Les ordres de grandeur d'équipe et
l'ordonnancement des briques sont une hypothèse de lecture, pas un plan.}

\end{document}
